%==============================================================================
%  Chapter 3 — Methodology
%
%  Structure:
%   1. Research design & document corpus
%   2. Analytical procedure (three-pass deductive → inductive → synthesis)
%   3. Use of generative AI in the research process
%   4. Quality, reflexivity and limitations
%==============================================================================
\chapter{Methodology}
\label{ch:methodology}

\section*{Chapter overview}

This chapter sets out how the thesis was done. It explains the research
design and the document corpus, the three-pass analytical procedure used
to read that corpus, the way a generative-AI assistant was used in the
workflow, and the quality criteria and limitations that frame what the
study can and cannot claim.

%==============================================================================
\section{Research design \& Document corpus}
\label{sec:meth-corpus}

This thesis takes a qualitative, interpretive approach. It treats policy
text as a primary object of analysis: the EIC Work Programme does not
simply describe the Pathfinder and STEP instruments, it constitutes them
by defining who counts as an applicant, what counts as innovation, and
what counts as success. The analytical task is therefore to read the
text closely and ask what assumptions and commitments it inscribes. This
fits the institutional-theory framework introduced in
Chapter~\ref{ch:background}, which treats categories such as
\enquote{innovation}, \enquote{market} and \enquote{state} as socially
constructed rather than given.

The primary subject of analysis during this thesis have been the EIC
Work Programme 2026, which functions to lay out the priorities and
direction for the different instruments managed by the EIC. The specific
instruments, subject of analysis, are scoped to the Pathfinder and STEP.
\\

The specific instruments, Pathfinder and STEP, were chosen to cover two
crucial inflection points;

\begin{enumerate}[leftmargin=*]
    \item The seed capital (Pathfinder)
    \item The scale-up pre-commercialisation (STEP)
\end{enumerate}

The two instruments sit at opposite ends of the EIC pipeline. Pathfinder
funds early-stage research where the technology is not yet proven; STEP
Scale Up funds companies that already have a working product and need
capital to scale. They share a common institutional home (the same
agency, the same Work Programme, the same governance structure), yet
they differ along several dimensions that the analysis can hold in
view: the type of applicant they address (research consortia versus
established companies), the form of capital they deploy (grant versus
blended grant-and-equity), the kind of evidence applicants are asked
to produce, and the evaluation criteria used to assess them. Holding
the institutional home constant while these dimensions vary is what
makes the pairing analytically useful: it allows the reading to attend
to how the textual treatment of an applicant, of evidence, and of
success shifts across the two instruments, without that variation
being confounded by a change of agency or governance regime. EIC
Accelerator and Transition are deliberately left out of the corpus
both to limit the scope and because they sit in the middle of the
pipeline, where the contrast that motivates the pairing would be
muted.

The study is positioned as an initial probe of the prior view stated
in Chapter~\ref{ch:introduction}. What the analysis can examine is
how the EIC inscribes the two instruments and what assumptions the
inscriptions carry: Pathfinder at the early stage at which the prior
view holds the dynamic to begin, and STEP at the scale-up stage to
which public capital has more recently been extended. The analysis
cannot, on its own, establish how those textual commitments shape
firm-level financing trajectories. Where the analysis surfaces
features that look consequential for the prior view, the appropriate
move is to mark them as warranting firm-level investigation. The
reflexive grounds for this framing are stated in
\cref{sec:meth-reflexivity}.

The corpus consists of a single document: the EIC Work Programme 2026.
The Work Programme implements two legal instruments, the STEP
Regulation (EU)~2024/795 and the Establishing Horizon Europe
Regulation (EU)~2021/695, and cites them as the binding basis for
the rules it sets out. These instruments are not themselves part of
the analysed material; they appear in the thesis only where the Work
Programme cites them, and the citation in those cases is to the Work
Programme. Secondary material such as think-tank commentary, press
coverage and EIC promotional material was also left out. The thesis
asks what the official text inscribes, not how it is received in the
wider field, and that scope follows from the question. The corpus and
the referenced legal instruments are listed in
Appendix~\ref{app:doc-corpus}.

The Work Programme was retrieved as the Commission's officially
published PDF (C(2025)\,7410) and kept locally as the canonical
reference. A plain-text working copy was generated to make search and
quoting easier; this is a mechanical step with no interpretive content,
but every passage quoted in Chapter~\ref{ch:results} was checked back
against the original PDF before being used.

The full project tree (corpus, codebooks, readthrough notes,
synthesis and the LaTeX report) is held in a public Git repository
\parencite{kirch2026thesisrepo}.
Substantive reorganisations are recorded in an append-only audit log
(\texttt{RESTRUCTURE\_LOG.md}). The numbered readthrough entries in
\texttt{Readthrough\_Notes\_v1.md} have stable identifiers, so any
finding reported in Chapter~\ref{ch:results} can be traced back to the
entry it comes from, and from there to the passage in the source PDF.
The chain of custody (PDF $\rightarrow$ extracted text $\rightarrow$
numbered entry $\rightarrow$ reported finding) is intentional and
gives the trustworthiness criteria discussed in
Section~\ref{sec:meth-quality} something concrete to point at.

%==============================================================================
\section{Analytical procedure}
\label{sec:meth-analysis}

The analysis ran in three sequential passes. The first applied a set of
framework codebooks deductively to the corpus and produced a numbered
set of coded entries. The second revisited the entries that had been
flagged as ambiguous, resolved them, and developed pattern-level
observations across the corpus. The third synthesised those patterns
into the named findings that Chapter~\ref{ch:results} presents and
Chapter~\ref{ch:discussion} interprets. Each pass fed the next, and the
whole procedure is recorded so any finding can be traced back to the
underlying passages.

\subsection{Pass 1: Deductive readthrough}
\label{sec:meth-pass1}

Pass~1 applied two codebooks to the corpus. The first covers Scott's
three pillars and DiMaggio and Powell's three types of isomorphism; the
second covers Thornton, Ocasio and Lounsbury's four institutional logics
(market, state, profession, science). Both codebooks are reproduced in
Appendix~\ref{app:codebook-deductive}.

The output of Pass~1 is a single working document,
\texttt{Readthrough\_Notes\_v1.md}, made up of numbered entries. Each
entry pairs a verbatim passage from the Work Programme with its codings
and a short note explaining the coding. The pass followed a conservative
standard: a logic was coded only where it clearly \emph{organised} the
passage, not merely where its vocabulary was present. The distinction
between organising logic and surface vocabulary was held throughout
the corpus, since the two do not always align. Passages that did not
fit cleanly were flagged with \texttt{[?]} markers and left for Pass~2
to resolve.

\paragraph{On identifying mimetic isomorphism in a document.}
The three isomorphism mechanisms are not textually visible in the same
way. \texttt{ISO\_COERCIVE} is anchored by regulatory citation and by
grant-agreement clauses; \texttt{ISO\_NORMATIVE} is anchored by named
evaluation communities, professional standards, or Commission-issued
guidance. \texttt{ISO\_MIMETIC} is different in kind: as
\textcite{dimaggio1983iron} specify, its diagnostic is structural
resemblance to a model in a reference field, not textual citation of
that model. Coding \texttt{ISO\_MIMETIC} from the corpus therefore
required a distinct identification strategy. Where the Work Programme
named an external referent, the code was straightforward. Where it did
not, the identification depended on matching the instrument's
structural features against models supplied by the secondary
literature (\cref{sec:bg-it-isomorphism}). The comparators used in
\cref{ch:results,ch:discussion} are analyst-supplied on this basis, and
the reading is explicit about the asymmetry throughout.

\subsection{Pass 2: Inductive update}
\label{sec:meth-pass2}

Pass~2 had two jobs. The first was to work through the \texttt{[?]}
flags from Pass~1: re-reading each ambiguous entry against both
codebooks and either resolving it to a coding or, where the ambiguity
was real, recording it as a productive analytical observation. The
second was to look across the resolved entries and develop
pattern-level observations: recurring textual moves that the
entry-by-entry coding could not see on its own
(\textcite{gioia2013seeking} provides the broad data-structure logic).

Pass~2 also produced seven \emph{coding rules} that discipline how the
v1 codes apply at the margin, without adding new codes. Five of them
underwrite the six patterns Chapter~\ref{ch:results} reports; the
other two discipline coding decisions without producing named patterns
of their own. The full set is reproduced in
Appendix~\ref{app:codebook-inductive}.

Inductive codes that emerged in Pass~2 (categories the deductive
codebooks did not anticipate) are documented in
Appendix~\ref{app:codebook-inductive}. The resolved flags and the
developed patterns are recorded in the same
\texttt{Readthrough\_Notes\_v1.md} file as Pass~1, so the deductive and
inductive layers sit alongside each other entry by entry.

\paragraph{Frozen v1 notes and logged consequences.}
The Pass~1 readthrough was frozen at the close of Pass~1 rather than
rewritten to reflect Pass~2 refinements. Freezing preserves the audit
trail: the reader can inspect what the deductive readthrough produced
without Pass~2 revisions overwriting the historical record. Two Pass~2
consequences would have implicated Pass~1 entries had the notes been
rewritten. [001]'s constructed-argument \texttt{LOGIC\_TENSION} coding
would be withdrawn under rule~1; [026] would attract
\texttt{ISO\_COERCIVE} co-coding under rule~3. Both are logged in
\texttt{Codebook\_v2\_Inductive.md}~\S9 (as v2-011 and v2-006
respectively) and are treated as in force for the analytical purposes
of Chapters~\ref{ch:results} and~\ref{ch:discussion}. Where
Chapter~\ref{ch:results} refers to a \enquote{frozen-notes
consequence}, this is what is meant.

\subsection{Pass 3: Synthesis and writing}
\label{sec:meth-pass3}

Pass~3 produced the cross-entry synthesis (\texttt{synthesis.md}) on
which Chapter~\ref{ch:results} draws. The synthesis tabulates the
logic profile by section of the corpus, groups recurring textual
moves into a small set of named patterns, and assembles a set of
STEP-specific observations alongside them. The patterns themselves,
their names, and the evidence on which each rests are presented in
Chapter~\ref{ch:results}; the interpretive readings built on top of
them are developed in Chapter~\ref{ch:discussion}.

Chapter~\ref{ch:discussion} draws on the synthesis rather than
re-reading the corpus, so the entry-level work feeds the chapter-level
interpretation through one mediating document.

%==============================================================================
\section{Use of generative AI in the research process}
\label{sec:meth-genai}

The thesis was written with the systematic support of a generative-AI
(GenAI) assistant.\footnote{Anthropic's \emph{Claude} family of large
language models (specifically Claude Opus 4.6 and Claude Opus 4.7),
accessed via the \emph{Claude Code} command-line interface. The
choice between the two versions was determined by the context-window
requirements of the specific task.} The
methodological literature on GenAI in qualitative research is still
unsettled, and supervisors and examiners reasonably want to know
exactly how the tool was used. This section sets out the position
taken, the tasks the tool was used for, how the analytical work was
kept anchored to the researcher's own judgement, and the limitations
that follow.

\subsection{Position}
\label{sec:meth-genai-position}

The published literature splits in roughly three directions.
\textcite{zhang2024redefining} and similar work see GenAI as a
productivity extension that, used carefully, can surface themes a
human coder might miss. \textcite{chatzichristos2025qualitative}
warns that the same move risks a quiet return to positivism: treating
coding as a repeatable mechanical operation hands back the reflexive,
interpretive work that defines qualitative analysis.
\textcite{nguyen2025engaged} go further and argue against putting any
AI between the researcher and the empirical material, on the grounds
that meaning is made in the slow work of reading and re-reading.

The position taken here is intermediate. The assistant is treated as
a tool of the researcher (closer in spirit to qualitative-analysis
software like MAXQDA or NVivo than to a co-analyst) and not as a
source of analytical judgement. The decisions that anchor the thesis
were made by the researcher: the institutional-theory readings of
Pathfinder and STEP, the identification and naming of cross-entry
patterns, and the interpretive framings developed in
Chapter~\ref{ch:discussion}. The tool was used for the work
\emph{around} those decisions: locating passages, keeping the codebook
consistent, formatting notes, drafting candidate prose for review, and
comparing across entries. Where Nguyen and Welch's warning is sharpest (handing
pattern identification over to the model), the workflow was set up so
that the model proposes and the researcher disposes.

\subsection{Tasks the tool was used for}
\label{sec:meth-genai-tasks}

The assistant was used for five categories of task:

\begin{enumerate}[leftmargin=*]
  \item \textbf{Codebook construction.} The deductive codebooks in
        \texttt{04\_codebook/output/} were drafted by the assistant from
        definitions the researcher supplied, then revised against the
        primary literature before being applied.
  \item \textbf{Corpus preparation.} The 188-page Work Programme PDF
        was converted to a plain-text working copy
        (\texttt{03\_corpus/output/WP2026\_extracted.txt}) to support
        search and quoting. No interpretive transformation was
        performed at this step.
  \item \textbf{Pass~1 deductive readthrough.} For each entry, the
        assistant produced a candidate coding against the framework,
        flagged ambiguous passages, and recorded emerging patterns.
        Every entry was reviewed, edited or rejected by the researcher
        before being kept.
  \item \textbf{Cross-entry synthesis.} Once Pass~1 covered enough of
        the corpus, the assistant was asked to identify recurring
        patterns across entries. These were treated as candidate
        observations and tested back against the underlying entries
        before being adopted, refined or dropped.
  \item \textbf{Drafting and editing.} Chapter prose, LaTeX scaffolding
        and bibliography management were drafted with assistance and
        then revised by the researcher. No claim, citation or finding
        appears in the final text without independent verification.
\end{enumerate}

The assistant was explicitly \emph{not} used to: generate or simulate
empirical material; keep coded entries that had not been reviewed by
the researcher; or formulate the thesis' central theoretical argument.

\subsection{Division of labour}
\label{sec:meth-genai-labour}

The workflow was structured around a deliberate division of labour
between researcher and tool. \Cref{tab:meth-ai-division} summarises the
allocation.

\begin{table}[ht]
  \centering
  \caption{Division of labour between researcher and GenAI assistant.}
  \label{tab:meth-ai-division}
  \begin{tabular}{p{0.45\linewidth} p{0.45\linewidth}}
    \toprule
    \textbf{Researcher} & \textbf{GenAI assistant} \\
    \midrule
    Research question and theoretical framing & Locating relevant passages in the corpus \\
    Selection of theoretical sources           & Drafting candidate codings against the codebook \\
    Final coding decisions                     & Maintaining schema and formatting consistency \\
    Identification and naming of patterns      & Cross-entry comparison and tabulation \\
    Interpretive argument                      & Drafting prose for researcher review \\
    Citation accuracy and source fidelity      & Bibliography formatting \\
    \bottomrule
  \end{tabular}
\end{table}

The split reflects the position above: the model accelerates work that
is structured, repetitive, or generative-of-candidates; the researcher
keeps every act of judgement on which the validity of the analysis
depends.

\subsection{Audit trail}
\label{sec:meth-genai-audit}

Three practices kept the AI-assisted workflow auditable.

First, the whole project tree is in Git, with the append-only
\texttt{RESTRUCTURE\_LOG.md} recording substantive reorganisations.
The numbered readthrough entries form a chronological record of every
coded passage; their stable numbering means any finding in
Chapter~\ref{ch:results} can be traced back to the entry it derives
from, and from there to the verbatim passage in the source PDF.

Second, candidate codings produced by the assistant were treated as
proposals and reviewed against the codebook. Where a suggestion was
ambiguous or did not fit cleanly, the entry was flagged with a
\texttt{[?]} marker and resolved during Pass~2. The resolved flags are
included as an appendix excerpt (Appendix~\ref{app:audit-trail}).

Third, no quotation, citation or factual claim in this thesis appears
without direct verification against the primary source. The full text
of the Work Programme is held locally and was the canonical reference
at every step. Secondary literature was read in the original rather
than via summary by the assistant. This is the standard mitigation for
the well-documented risk of fabricated citations and factual drift in
large-language-model outputs \parencite{nguyen2025engaged}.

\subsection{Limitations of AI assistance}
\label{sec:meth-genai-limits}

Three limitations are acknowledged.

The first is the \emph{positivist drift} warned about by
\textcite{chatzichristos2025qualitative}. AI-assisted coding can
quietly re-import the assumption that meaning lives in the text and
is recoverable by repeatable procedure. The workflow tries to
mitigate this by keeping all interpretive moves with the researcher
and by working within an institutional-logics framework that is itself
explicit about the social construction of categories. The risk is
reduced, not eliminated.

The second is \emph{anchoring}. Even where the researcher revises a
model-generated coding, the initial suggestion may bias the direction
of attention. Pass~2 was introduced in part to counter this, by
re-reading entries against unresolved questions rather than against
the original codings.

The third is \emph{opacity}. The model's reasoning is not directly
inspectable and the conditions of its training are only partially
known. For that reason the assistant is treated as a fallible
drafting collaborator whose outputs require verification, not as an
independent source of evidence.

This account is intended to satisfy the responsible-use expectations
that are now common across the qualitative-research community: a
transparent disclosure of GenAI use, a clear statement of researcher
responsibility for the final text, and an explanation of how risks
to integrity and interpretation were managed.

%==============================================================================
\section{Quality, reflexivity and limitations}
\label{sec:meth-quality}

\subsection{Quality criteria}
\label{sec:meth-quality-criteria}

Trustworthiness in this study is addressed through Lincoln and Guba's
four standard criteria for qualitative research \parencite{lincoln1985naturalistic}:
credibility, transferability, dependability and confirmability.

\emph{Credibility} (that the analysis stays faithful to what the
text actually says) was supported in two ways. The numbered entry
format pairs each coding with its verbatim passage, so any coding can
always be checked against its source. As an additional check, the
researcher also worked through the Work Programme by hand on a
printed copy with coloured markers, and that manual coding is
cross-referenced against the digital coding in
\texttt{colour-codes-cross-reference.md}. Where the two passes agreed,
the coding is more confident; where they diverged, the divergence is
documented and explained.

\emph{Transferability} (that the reader can judge whether the
findings travel to other settings) is supported by thick
description and verbatim quotation throughout Chapter~\ref{ch:results},
and by an explicit scope statement of where the findings might and
might not apply (Chapter~\ref{ch:discussion}, Boundary conditions).

\emph{Dependability} (that the analytical procedure is stable and
could in principle be followed by someone else) rests on the
codebook (Appendix~\ref{app:codebook-deductive}), the audit trail in
\texttt{Readthrough\_Notes\_v1.md}, and the append-only
\texttt{RESTRUCTURE\_LOG.md}. The procedure is documented; the inputs
and outputs of each stage are versioned.

\emph{Confirmability} (that the findings are grounded in the data
and not in the researcher's prior commitments) is addressed by the
entry-level coding rationales, the manual cross-reference as an
independent check, and the explicit division of labour with the AI
assistant set out in \cref{sec:meth-genai-labour}.

\subsection{Reflexivity}
\label{sec:meth-reflexivity}

My background shaped the choice of object and the framing of the
prior view stated in Chapter~\ref{ch:introduction}. More than four
years of work in venture building, investing and consulting shaped
the way I came to read the European public-funding layer, and the
Pathfinder and STEP instruments were chosen for analysis with that
background in view.

The reading the thesis lands on in Chapter~\ref{ch:discussion}
points in the same direction as the prior view that brought me to
the topic. The match is not in itself evidence of bias, but it is
the obvious risk to flag. Two safeguards were used. First, the
Chapter~\ref{ch:discussion} argument is tested explicitly against
rival readings (\cref{sec:disc-rival}). Second, the coding work was
anchored in what the Work Programme text inscribes rather than in the
prior view: every coded entry pairs its coding with a verbatim
passage, and the hand-coded manual cross-reference
(\cref{sec:meth-quality-criteria}) gives the reader an independent
check on whether the textual evidence supports the readings.


\subsection{Ethics}
\label{sec:meth-ethics}

A documentary design on publicly available EU policy texts avoids the
human-subjects ethics concerns that come with interview research. The
only responsible-use question that remains concerns the use of
generative AI, which is addressed in \cref{sec:meth-genai}.

\subsection{Methodological limitations}
\label{sec:meth-limits}

Four methodological limitations are acknowledged here. Substantive
limitations of the findings themselves are kept for
Chapter~\ref{ch:conclusion}.

\begin{enumerate}[leftmargin=*]
  \item \textbf{Single-source scope.} The corpus is a single document.
        The two regulations cited within it were treated as referenced
        legal background, not as coded material. The study sees what
        the Work Programme says, not how it is enacted on the ground.
  \item \textbf{Moving target.} STEP only recently became operational.
        Some of the textual moves analysed here may look different
        once a few funding cycles have run and the instrument has been
        revised.
  \item \textbf{English-language only.} The Work Programme is published
        in English; this study works only from the English text and
        does not compare across language versions.
  \item \textbf{Single-coder design.} The analysis was carried out by
        one researcher. The manual cross-reference described under
        \emph{credibility} is the partial compensation for the absence
        of a second coder; it does not replace inter-rater reliability
        in the strict sense.
\end{enumerate}
